\scanpage{73}
\mspage{II.21}
% archive sheet number: 72
% lines 1–5
2.5. Dans ce qui suit, nous supposerons $\dim X=1$, donc $\dim X'=1$. Quitte à faire encore une extension finie sur $\unsure{k|h}$, on peut supposer les points $x_i$ ($1\leq i\leq n$) de $X'-X$ rationnels sur $k$. On a alors une suite exacte bien connue
% lines 6–7
\noindent (2.5.1)
\diagram{\begin{tikzcd}[ampersand replacement=\&,column sep=small]
\unsure{0|1} \arrow[r,dash] \& H^1(\bar X',\unsure{\mathbb Z_\ell|Z_\ell}) \arrow[r] \& H^1(\bar X,\mathbb Z_\ell) \arrow[r] \& \displaystyle\prod_{\deleted{\unsure{0}}1\leq i\leq n}H_x^2(\bar X,\mathbb Z_\ell) \arrow[r] \& H^2(\bar X',\mathbb Z_\ell) \arrow[d] \\
\&\&\&\&0
\end{tikzcd}}
% lines 8–10
\deleted{qui est} \inserted{dont les flèches sont} évidemment compatibles avec l'action de \deleted{$\unsure{\pi_0}$} \deleted{$\unsure{k}$} $\pi_0=\operatorname{Gal}(\bar k/k)$. De plus, on a un isomorphisme canonique
% line 11
\[
H_x^2(\bar X,\mathbb Z_\ell)\simeq\mathbb Z_\ell(-1)\qquad \deleted{\unsure{(\varprojlim_\nu{}_{\ell^\nu}\bar k^*)}}\tag*{(2.5.2)}
\]
% line 12
où on a
% line 13
\[
\deleted{\unsure{H^1}}\mathbb Z_\ell(-1)=\mathbb Z_\ell(1)^{\otimes(-1)},\qquad \mathbb Z_\ell(1)=\varprojlim_{\unsure{\nu|v}}({}_{\ell^\nu}\bar k^{\unsure{*|\times}}).\tag*{(2.5.3)}
\]
% lines 14–19
(De plus, on a un isom. canonique $H^2(\bar X',\mathbb Z_\ell)\simeq\mathbb Z_\ell(-1)$, et l'hom $\prod\longrightarrow H^{\unsure{1|2}}(\bar X',\mathbb Z_\ell)$ s'identifie à l'hom somme). L'isomorphisme (2.5.2) est encore compatible avec l'action de $\pi_0$, opérant sur le deuxième m\textsuperscript{bre} de la façon explicitée par la formule (2.5.3). \deleted{\unsure{Il s'}}

% lines 20–23
Il s'ensuit qu'en tant que $\overset{\deleted{\unsure{-}}}{\pi}_0$-module $\ell$-adique, $H^1(\bar X,\mathbb Z_\ell)$ est une extension de $\mathbb Z_\ell(1)^{n-1}$ par $H^1(\bar X',\mathbb Z_\ell)$. Comme l'opération de $\pi_0$ sur $\mathbb Z_\ell(1)$ est non triviale, il s'ensuit \inserted{déjà} qu'on a
